
\section*{\SGAInlineRefTarget{sga3:i:section:3}3. The category of \(\mathbf O\)-modules, the category of
\(\mathbf G\)-\(\mathbf O\)-modules}
\addcontentsline{toc}{section}{3. The category of O-modules, the category of G-O-modules}

\medskip
\noindent\SGAInlineRefTarget{sga3:i:definition:3-1}\textbf{Definition 3.1.}
Let \(\mathbf O\) and \(\mathbf F\) be two objects of
\(\widehat{\mathscr C}\). We say that \(\mathbf F\) is a
\(\widehat{\mathscr C}\)-module over the
\(\widehat{\mathscr C}\)-ring \(\mathbf O\), or, for short, an
\(\mathbf O\)-module, if, for every
\(S\in\operatorname{Ob}(\mathscr C)\), the set \(\mathbf O(S)\) has been
endowed with a ring structure and \(\mathbf F(S)\) with a module structure
over that ring, in such a way that for every arrow
\(S'\to S''\) in \(\mathscr C\), the map
\(\mathbf O(S'')\to\mathbf O(S')\) is a ring homomorphism and
\(\mathbf F(S'')\to\mathbf F(S')\) is a homomorphism of abelian groups
compatible with that ring homomorphism.

Once \(\mathbf O\) is fixed, one defines in the usual way a morphism between
two \(\mathbf O\)-modules \(\mathbf F\) and \(\mathbf F'\), hence the
abelian group
\(\operatorname{Hom}_{\mathbf O}(\mathbf F,\mathbf F')\), and the category
of \(\mathbf O\)-modules, denoted \((\mathbf O\text{-Mod.})\).
Note that if \(\mathbf O_0\) is the \(\widehat{\mathscr C}\)-ring defined by
\[
  \mathbf O_0(S)=\mathbb Z
\]
(which must not be confused with the functor associated with the constant
object \(\mathbb Z\)), then the category of \(\mathbf O_0\)-modules is
isomorphic to the category of commutative
\(\widehat{\mathscr C}\)-groups. The category \((\mathbf O\text{-Mod.})\) is
abelian.

\par\medskip\noindent\SGAInlineRefTarget{sga3:i:lemma:3-1-1}%
\emph{Addition} (after the Gille--Polo re-edition (2011), Exposé~I, 3.1.1
and 3.1.2; our formulation). Kernels, cokernels, images and arbitrary
direct sums of \(\mathbf O\)-modules can be formed pointwise
(\(S\mapsto\operatorname{Ker}(u(S))\), and so on); they have the required
universal properties, a morphism of \(\mathbf O\)-modules being a family of
module homomorphisms compatible with the maps
\(\mathbf F(S'')\to\mathbf F(S')\). The canonical morphism from coimage to
image is an isomorphism, being one pointwise; so
\((\mathbf O\text{-Mod.})\) is abelian, a sequence being exact if and only
if it is exact at every \(S\). As filtered direct limits of modules are
exact, it satisfies axiom \((\mathrm{AB}\,5)\) of A.~Grothendieck, \textit{Sur quelques points d'algèbre homologique}, Tôhoku Math.\ J.\ \textbf{9} (1957), 119--221, 1.5. If
\(\mathscr C\) is small, it has a generator. For
\(S\in\operatorname{Ob}(\mathscr C)\), let \(\mathbf O_{(S)}\) be the
\(\mathbf O\)-module \(S'\mapsto\mathbf O(S')^{(\operatorname{Hom}(S',S))}\),
free on \(\mathbf h_S\); by 1.1, the morphisms
\(\mathbf O_{(S)}\to\mathbf F\) correspond to the elements
\(x\in\mathbf F(S)\), the basis element \(\operatorname{id}_S\) going to
\(x\). The resulting morphism
\(\bigoplus_{S,\,x\in\mathbf F(S)}\mathbf O_{(S)}\to\mathbf F\) is surjective
at every \(S\); so \(\bigoplus_S\mathbf O_{(S)}\) is a generator, and
\((\mathbf O\text{-Mod.})\) has enough injective objects
(\textit{loc.\ cit.}, 1.10.1).
\par\medskip

\medskip
\noindent\SGAInlineRefTarget{sga3:i:definition:3-1-4}Notice that if \(\mathbf F\) is an \(\mathbf O\)-module, then for every
\(S\in\operatorname{Ob}(\mathscr C)\), \(\mathbf F_S\) is an
\(\mathbf O_S\)-module. We may therefore define an abelian
\(\widehat{\mathscr C}\)-group
\(\underline{\operatorname{Hom}}_{\mathbf O}(\mathbf F,\mathbf F')\) by
\[
  \underline{\operatorname{Hom}}_{\mathbf O}(\mathbf F,\mathbf F')(S)
    =\operatorname{Hom}_{\mathbf O_S}(\mathbf F_S,\mathbf F'_S).
\]
In the same way one defines the objects
\[
  \underline{\operatorname{Isom}}_{\mathbf O}(\mathbf F,\mathbf F'),
  \qquad
  \underline{\operatorname{End}}_{\mathbf O}(\mathbf F),
  \qquad\text{and}\qquad
  \underline{\operatorname{Aut}}_{\mathbf O}(\mathbf F),
\]
the last of which is a \(\widehat{\mathscr C}\)-group under the structure
induced by composition of automorphisms.

\medskip
\noindent\SGAInlineRefTarget{sga3:i:definition:3-2}\textbf{Definition 3.2.}
Let \(\mathbf O\) be a \(\widehat{\mathscr C}\)-ring,
\(\mathbf F\) an \(\mathbf O\)-module, and \(\mathbf G\) a
\(\widehat{\mathscr C}\)-group. A \(\mathbf G\)-\(\mathbf O\)-module
structure on \(\mathbf F\) is a \(\mathbf G\)-object structure such that,
for every \(S\in\operatorname{Ob}(\mathscr C)\) and every
\(g\in\mathbf G(S)\), the automorphism of \(\mathbf F(S)\) defined by
\(g\) is an automorphism of its \(\mathbf O(S)\)-module structure.

Equivalently, the morphism of \(\widehat{\mathscr C}\)-groups
\[
  \rho:\mathbf G\longrightarrow
    \underline{\operatorname{Aut}}(\mathbf F)
\]
defined in \SGARefLink{sga3:i:subsection:2-3}{2.3} maps \(\mathbf G\) into the
\(\widehat{\mathscr C}\)-subgroup
\(\underline{\operatorname{Aut}}_{\mathbf O}(\mathbf F)\) of
\(\underline{\operatorname{Aut}}(\mathbf F)\). Thus, giving a
\(\mathbf G\)-\(\mathbf O\)-module structure on the
\(\mathbf O\)-module \(\mathbf F\) is the same as giving a morphism of
\(\widehat{\mathscr C}\)-groups
\[
  \rho:\mathbf G\longrightarrow
    \underline{\operatorname{Aut}}_{\mathbf O}(\mathbf F).
\]
One defines naturally the abelian group
\(\operatorname{Hom}_{\mathbf G\text{-}\mathbf O}
  (\mathbf F,\mathbf F')\), and hence the additive category of
\(\mathbf G\)-\(\mathbf O\)-modules, denoted
\((\mathbf G\text{-}\mathbf O\text{-Mod.})\). The reader will notice that the latter
may also be defined as the category of \(\mathbf O[\mathbf G]\)-modules,
where \(\mathbf O[\mathbf G]\) is the algebra of the
\(\widehat{\mathscr C}\)-group \(\mathbf G\) over the
\(\widehat{\mathscr C}\)-ring \(\mathbf O\), whose definition is clear.\footnote{Translator's note: the Gille--Polo re-edition (2011), Exposé~I, notes
here that \((\mathbf G\text{-}\mathbf O\text{-Mod.})\), being the category
of \(\mathbf O[\mathbf G]\)-modules, is therefore abelian, satisfies
\((\mathrm{AB}\,5)\), and has enough injective objects if \(\mathscr C\) is
small (Addition in 3.1, which applies to modules over any
\(\widehat{\mathscr C}\)-ring, commutative or not); this is used in 5.2.}

All the constructions of this section specialize immediately to the case
where \(\mathbf G\) (or \(\mathbf O\), or both) is representable by an
object of \(\mathscr C\), which is thereby endowed with the corresponding
algebraic structure.

We have treated briefly the principal algebraic structures that occur later
in this seminar. For the others (the structure of an \(\mathbf O\)-Lie
algebra, for example), we believe that the reader has now seen enough
examples in this section to make the general mechanism sketched in \SGARefLink{sga3:i:subsection:2-2}{2.2} work
in each particular case.

We shall now apply what we have just done to the category of schemes, denoted
\((\mathbf{Sch})\), and more generally to the categories
\((\mathbf{Sch})_{/S}\).

